\chapter{Intravenous Cannulation}

\section*{Overview}
Intravenous (IV) cannulation is the insertion of a plastic cannula into a peripheral vein to provide access to the bloodstream. It allows the administration of fluids, medications, blood products, and contrast media, and is a fundamental skill in acute and hospital care.
\section*{Alternative Names}
Peripheral intravenous line insertion; peripheral IV access; PIVC insertion; IV start
\section*{Principle}
A hollow needle with an overlying plastic catheter is advanced into a distended vein under tourniquet control. Once a flashback of blood confirms entry into the lumen, the needle is withdrawn and the catheter is advanced fully into the vein, where it is secured and connected to infusion tubing.
\section*{History}
The practice of intravenous infusion began with early experimental injections in the 17th century and became a routine therapy after the introduction of the indwelling intravenous cannula and improved equipment in the mid-20th century.
\section*{Why It Is Done}
Ordered to deliver intravenous fluids and electrolytes, administer medications and antibiotics, give blood products, provide parenteral nutrition, administer contrast for imaging, and establish emergency vascular access when oral or other routes are unsuitable.
\section*{Is This a Primary or Secondary Test or Procedure}
A primary supportive procedure that establishes a route for therapy; it is an access procedure rather than a diagnostic or curative intervention.
\section*{How It Is Performed}
A tourniquet is applied and a vein on the dorsal hand, forearm, or antecubital fossa is selected. The site is disinfected and the skin is punctured with the needle-catheter assembly bevel up at a shallow angle. On observing flashback, the needle is removed, the catheter is threaded into the vein, the tourniquet is released, and the catheter is flushed, capped, and secured with a transparent dressing.
\section*{Preparation}
Informed consent and verification of the planned therapy. The operator selects the smallest adequate gauge and warms the limb or uses a tourniquet to improve vein distension when needed.
\section*{Contraindications and Precautions}
Avoid cannulating a limb with an arteriovenous fistula, lymphedema, thrombophlebitis, cellulitis, or previous mastectomy on that side. Use caution in fragile or rolling veins, in anticoagulated patients, and in children, and do not persist with repeated failed attempts.
\section*{Risks}
Pain, bleeding, hematoma, phlebitis, extravasation of fluid or medication, infection, thrombophlebitis, and rarely nerve or arterial injury. Cannula occlusion, dislodgement, and infiltration are common complications during therapy.
\section*{Aftercare and Recovery}
The line is labelled with the insertion date and the site is inspected regularly for signs of phlebitis or infection. Peripheral cannulas are typically replaced or removed after 72-96 hours or when no longer needed.
\section*{Outcome and Follow-up}
Cannulation success is confirmed by blood flashback, free flush, and secure placement; continued patency is indicated by good flow and the absence of pain or swelling during infusion.
\section*{Expected Results}
A patent cannula that flushes freely without pain, swelling, or leakage, and an insertion site that remains clean with no erythema, tenderness, or discharge.
\section*{Follow-up}
The site is monitored for complications, and cannula removal with site care is performed when therapy is complete. Persistent or difficult access may require ultrasound guidance or central venous catheter placement.
\section*{When to Seek Medical Care}
Follow the procedure team's specific instructions. Seek urgent medical assessment for severe or worsening pain, uncontrolled bleeding, fever or signs of infection, trouble breathing, chest pain, fainting, new weakness or confusion, inability to urinate, or any rapidly worsening symptom. The appropriate warning signs depend on the procedure and the patient's medical condition.

\section*{References}
\begin{enumerate}
  \item MedlinePlus. Intravenous line. U.S. National Library of Medicine; 2025.
  \item Current Medical Diagnosis and Treatment. McGraw-Hill; 2025.
\end{enumerate}

\clearpage
